\documentclass[11pt]{article}
\usepackage[a4paper,margin=27mm]{geometry}
\usepackage{amsmath,amssymb,mathtools}
\usepackage{fontspec,unicode-math}
\setmainfont{Latin Modern Roman}
\setmathfont{Latin Modern Math}
\usepackage{xurl}
\usepackage[hidelinks]{hyperref}
\newcommand{\sref}[2]{\hyperlink{source:#1:#2}{[S#2]}}
\newcommand{\cataloguescope}[1]{\paragraph{Mathematical question.}#1}
\hypersetup{pdftitle={Problem statement}}
\begin{document}
\section{Are there infinitely many cluster primes}\label{statement}

\subsection{Definitions and mathematical statement}

\paragraph{Shared notation.}$\mathbb N=\{1,2,\ldots\}$, and $\mathbb Z,\mathbb Q,\mathbb R,\mathbb C$ denote the integers, rationals, reals and complex numbers. Any other conventions are specified locally.


An odd prime $p$ is a cluster prime if, for every even positive integer $d$ with $d\le p-3$, there exist primes $q,r\le p$ with $q-r=d$. The prime $2$ is excluded by convention; for $p=3$ the condition is empty.
\paragraph{Statement recorded in the supplied source.}
Are there infinitely many cluster primes? Equivalently, is the set of odd primes satisfying this universal difference condition unbounded? \sref{statement}{2}
\subsection{Short English statement}
Do infinitely many odd primes have the property that every possible positive even difference up to three less than the prime is realized by two primes no larger than it?
\subsection{Sources}
\begin{description}
\item[\hypertarget{source:statement:1}{S1}] ULAM.AI and the UnsolvedMath Contributors, UnsolvedMath: A Curated Collection of Open Mathematics Problems, catalogue entry. CC BY 4.0. \url{https://huggingface.co/datasets/ulamai/UnsolvedMath}.
\item[\hypertarget{source:statement:2}{S2}] Wikipedia contributors, \emph{Cluster prime}, revision 1377809312 (2026) Lead and infinitude question. \url{https://en.wikipedia.org/w/index.php?title=Cluster+prime&oldid=1377809312}.
\end{description}
\label{end:statement}
\end{document}
